\subsection{Replication Ring}
\label{subsec:replication-ring}

To ensure database availability, each key needs to be replicated across
multiple nodes while satisfying the following requirements:

\begin{enumerate}
    \item Each key is replicated on $N$ nodes, each from a distinct
        availability zone.
    \item $w_i \leq s_i$ for each node $i$, where $w_i$ is the workload on node
        $i$ and $s_i$ is the bandwidth allocated by that node for the table,
        both measured in requests per second.
\end{enumerate}

The workload on the $i$-th node is $w_i = \sum_{k \in K_i} p_k \cdot B$ where
$K_i$ is the set of keys assigned to node $i$, $p_k$ is the probability that
key $k$ is requested and $B$ is the total bandwidth of the table, in requests
per second.

An adverse distribution of keys in requests would make it impossible to scale
the table on multiple nodes, for example if node $i$ is assigned a key $k$ such
that $p_k = 1$ we get that $w_i = B \leq s_i$, requiring that node $i$ is able
to handle the entire table workload.
To avoid this problem, we assume that keys are uniformly distributed, i.e.,
$p_k = \frac{1}{|K|}$ for each key $k$ in the table, where $K$ is set of keys
in the table.

Under this assumption, the workload on node $i$ is $w_i = \frac{|K_i|}{|K|}
\cdot B$ so we can rewrite the second requirement as $\frac{|K_i|}{|K|} \leq
\frac{s_i}{B}$.

Since we don't want to waste resources, we are going to strengthen the
requirement and enforce that $\frac{|K_i|}{|K|} = \frac{s_i}{B}$.
Summing over all nodes from the same availability zone $z$, we get $\sum_{i :
\mathrm{az}(i) = z} \frac{|K_i|}{|K|} = \sum_{i : \mathrm{az}(i) = z}
\frac{s_i}{B}$, since each key is assigned to at most one node from the same
availability zone, we have $\sum_{i : \mathrm{az}(i) = z} |K_i| \leq |K|$,
hence $\sum_{i : \mathrm{az}(i) = z} s_i \leq B$, and $s_i \leq B$ for each
node $i$.

%\begin{equation}
%\label{eq:total-bandwidth}
%S = B \cdot N
%\end{equation}

%\begin{lemma}[AZ feasibility]
%    \label{lem:az-feasibility}
%    Under the two requirements above, for each availability zone $z$ the
%    following condition must hold:
%    \[
%        \sum_{i : \mathrm{az}(i) = z} s_i \leq b
%    \]
%\end{lemma}
%\begin{proof}
%    Summing the second requirement over all nodes from the same availability
%    zone $z$, we get $\sum_{i : \mathrm{az}(i) = z}
%    \frac{|K_i|}{|K|} = \sum_{i : \mathrm{az}(i) = z} \frac{s_i}{B}$. Since
%    each key is assigned to at most one node from the same availability zone,
%    we have $\sum_{i : \mathrm{az}(i) = z} |K_i| \leq |K|$, hence
%    $\sum_{i : \mathrm{az}(i) = z} s_i \leq b$.
%\end{proof}

Let's define $S = \sum s_i$, the total bandwidth allocated for the table across
all nodes.
By summing the second requirement over all nodes, we obtain $\frac{1}{|K|} \sum
|K_i| = \frac{S}{B}$.
Since each key is assigned to exactly $N$ nodes, we have $\sum |K_i| = N
\cdot |K|$, and so $S = B \cdot N$.

To solve this problem we will use the \emph{replication ring}, this structure
will require that values $S$, $B$ and all $s_i$ are integers with a granularity
of $1$ request per second.
We define the \emph{replication ring} as a circular array of size $S$ where for
each position in the ring we assign a node $i$ such that:

\begin{enumerate}
    \item Every consecutive interval of length $N$ in the ring contains
        distinct nodes with different availability zones.
    \item Each node $i$ occupies exactly $s_i$ positions in the ring.
\end{enumerate}

Given a key $k$, we find its starting position $p$ in the ring using its hash
value (using \texttt{XXH3}\footnote{\url{https://xxhash.com}}, a fast
non-cryptographic hash algorithm) and taking the result modulo $S$.
The key is then assigned to the $N$ nodes immediately following the
computed position, i.e., to the nodes at positions $p, p + 1, \dots, p + n - 1$
(wrapping around the ring if necessary).

Since each interval of length $l$ contains distinct nodes, each key is assigned
to $N$ nodes with different availability zones, satisfying the first
requirement.
Furthermore, the probability that a key is assigned to position $p$ in the ring
is $\frac{l}{S} = \frac{1}{B}$, so the expected number of keys assigned to the
position is $|K| \cdot \frac{1}{B}$.
Since each node $i$ occupies exactly $s_i$ positions in the ring and each key
is assigned at most once to the same node, the expected number of keys assigned
to node $i$ is $|K_i| = |K| \cdot \frac{s_i}{B}$, thus satisfying the second
requirement.

\subsubsection{Explicit construction}
\label{subsubsec:explicit-construction}

First we are going to order the nodes, we will sort them such that nodes from
the same availability zone are grouped together and such that availability
zones $z$ with $\sum_{i : \mathrm{az}(i) = z} s_i = B$ are at the beginning of
the order.

Then nodes are assigned to the ring iteratively in the order defined above.
Each node is assigned to $s_i$ positions, proceeding from the last assigned
position to the previous node (or from an arbitrary position for the first
node) and occupying one position every $N$.
If that position is already taken by another node, we skip to the next; it is
guaranteed that it will never be necessary to skip two consecutive positions.

Since each availability zone $z$ has bandwidth $\sum_{i : \mathrm{az}(i) = z}
s_i \leq B$, and the ring has size $S = B \cdot N$, it is guaranteed that the
nodes from the same availability zone $z$ are assigned to $\sum_{i :
\mathrm{az}(i) = z} s_i$ positions without completing a full cycle of the ring,
so all positions occupied by nodes of the same availability zone are at least
$N$ positions apart.
Therefore, this construction satisfies the properties described above for the
\emph{replication ring}.

\noindent In pseudocode, the construction of the ring can be represented as in Listing~\ref{lst:explicit-ring}.

\begin{lstlisting}[
    language=Python,
    float=htbp,
    caption={Explicit ring construction},
    label={lst:explicit-ring}
]
def build_ring(table):
  S = table.bandwidth * table.n
  ring = [None] * S
  curr_pos = 0

  # Assume nodes are already sorted
  for i in range(table.m):
    for _ in range(table.s[i]):
      if ring[curr_pos] is not None:
        curr_pos = (curr_pos + 1) % S
      assert ring[curr_pos] is None
      ring[curr_pos] = i
      curr_pos = (curr_pos + table.n) % S

  return ring
\end{lstlisting}

This construction requires $O(S)$ time and space.

\paragraph{Example}
Let's see an example of ring construction for a table with parameters:
\[
    N = 3,\quad
    b = 4,\quad
    S = B \cdot N = 12,
\]
and five nodes all from distinct availability zones with bandwidth allocated
for the table:
\[
    s_0 = 4,\; s_1 = 3,\; s_2 = 2,\; s_3 = 1,\; s_4 = 2.
\]

We start with a circular array $r_0, r_1, \dots, r_{11}$ of length $S = 12$.
Each node $i$ is assigned to $s_i$ positions, jumping $N = 3$ slots from the
last assigned position.

\begin{itemize}
    \item Node $0$ is assigned to $s_0 = 4$ positions: $r_0$, $r_{0+3} = r_3$,
        $r_{3+3} = r_6$, and $r_{6+3} = r_9$.
    \item Node $1$ is assigned to $s_1 = 3$ positions: $r_{9+3} = r_{0}$ is
        already taken, so we skip to $r_1$, then $r_{1+3} = r_4$, and
        $r_{4+3} = r_7$.
    \item Node $2$ is assigned to $s_2 = 2$ positions: $r_{7+3} = r_{10}$ and
        $r_{10+3} = r_1$ is already taken, so we skip to $r_2$.
    \item Node $3$ is assigned to $s_3 = 1$ position: $r_{2+3} = r_5$.
    \item Node $4$ is assigned to $s_4 = 2$ positions: $r_{5+3} = r_8$ and
        $r_{8+3} = r_{11}$.
\end{itemize}

%\[
%    \begin{array}{rl}
%        \text{Node }0\;(s_0=4): & r_{0\phantom{0}} = 0 \\
%                                & r_{3\phantom{0}} = 0 \\
%                                & r_{6\phantom{0}} = 0 \\
%                                & r_{9\phantom{0}} = 0 \\[0.5ex]
%        \text{Node }1\;(s_1=3): & r_{1\phantom{0}} = 1 \\
%                                & r_{4\phantom{0}} = 1 \\
%                                & r_{7\phantom{0}} = 1 \\[0.5ex]
%        \text{Node }2\;(s_2=2): & r_{10} = 2 \\
%                                & r_{2\phantom{0}} = 2 \\[0.5ex]
%        \text{Node }3\;(s_3=1): & r_{5\phantom{0}} = 3 \\[0.5ex]
%        \text{Node }4\;(s_4=2): & r_{8\phantom{0}} = 4 \\
%                                & r_{11} = 4
%    \end{array}
%\]

Finally, we obtain:
\[
    \begin{array}{c|cccccccccccc}
        p   & 0 & 1 & 2 & 3 & 4 & 5 & 6 & 7 & 8 & 9 & 10 & 11 \\ \hline
        r_p & 0 & 1 & 2 & 0 & 1 & 3 & 0 & 1 & 4 & 0 & 2  & 4
    \end{array}
\]

It is easy to verify that every window of length $N = 3$ contains distinct
nodes and that each node $i$ occupies exactly $s_i$ positions.

\subsubsection{Implicit construction}
\label{subsubsec:implicit-construction}

Since the allocated bandwidth can be considerably high, constructing the ring
explicitly may be really expensive.
To avoid this, we look for an algorithm that directly computes the node
assigned to a position $p$ in the ring.
Let $o_p$ be the iteration at which the $p$-th position is assigned.
We observe that the explicit construction assigns nodes in the following order:
first the multiples of $N$ in increasing order, then the numbers congruent to
$1$ modulo $N$ in increasing order, then the numbers congruent to $2$ modulo
$N$, and so on up to $N-1$.
It follows that $o_p = \left\lfloor \frac{p}{N} \right\rfloor + (p \bmod N)
\cdot N$.
Now we need to find the $o_p$-th assigned node, i.e., the node $i$ such that
$\sum_{j=0}^{i-1} s_j \leq o_p < \sum_{j=0}^{i} s_j$.

In pseudocode, the algorithm to find the node assigned to a position $j$ in the
ring can be represented as in Listing~\ref{lst:implicit-ring}.

\begin{lstlisting}[
    language=Python,
    float=htbp,
    caption={Implicit ring construction},
    label={lst:implicit-ring}
]
def get_node_at_position(table, j):
    n = table.n
    b = table.bandwidth
    o_p = floor(p / n) + (p % n) * b

    sum = 0
    for node in range(table.m):
        sum += table.s[node]
        if sum > o_p:
            return node
\end{lstlisting}

We can therefore compute the node assigned to a position $p$ in the ring in
time $O(m)$ and constant memory, without explicitly building the ring.

The algorithm can be further optimized using a binary search on the prefix sums
of $s_i$, assuming they are already computed and stored.
In this case the algorithm runs $O(\log m)$ time and $O(m)$ memory.

\paragraph{Example}
Let's apply the implicit construction to the example in the previous section.
We start by computing the prefix sums $c_i = \sum_{j=0}^{i-1}
s_i$, obtaining
\[
    \begin{array}{c|cccccc}
        i   & 0 & 1 & 2 & 3 &  4 &  5 \\ \hline
        s_i & 4 & 3 & 2 & 1 &  2 &    \\ \hline
        c_i & 0 & 4 & 7 & 9 & 10 & 12
    \end{array}
\]

We then compute
\[
    o_p = \Bigl\lfloor \frac{p}{N} \Bigr\rfloor + (p \bmod N) \cdot B.
\]
For example, for $p = 5$ we have
\[
    o_5 = \Bigl\lfloor \tfrac{5}{3} \Bigr\rfloor + (5 \bmod 3) \cdot 4 =
    1 + 2\cdot4 = 9.
\]

We then look for $i$ such that $c_i \leq o_5 < c_{i+1}$, finding $c_3 = 9 \leq
9 < 10 = c_4$, hence the node at $p = 5$ is node~$3$.

In this way we can compute $r_p$ for any $p$ without paying the cost of
building the entire ring.

Repeating the algorithm for all $p = 0, \dots, 11$ we obtain:
\[
    \begin{array}{c|cccccccccccc}
        p   & 0 & 1 & 2 & 3 & 4 & 5 & 6 & 7 &  8 & 9 & 10 & 11 \\ \hline
        o_p & 0 & 4 & 8 & 1 & 5 & 9 & 2 & 6 & 10 & 3 &  7 & 11 \\ \hline
        r_p & 0 & 1 & 2 & 0 & 1 & 3 & 0 & 1 &  4 & 0 &  2 &  4
    \end{array}
\]

This matches the result of the explicit construction.
